%!TEX root = ../../main.tex
\section{시간 계약과 누출 방지}
\label{sec:data-contract}

예측 시점은 $\tcut = t_{\text{첫 킬}} - B$, $B = \constB$초이며, $q$와 승률·시간 기준선은 $\tcut$ 또는 그 직전의 상태 스냅숏 $\Spre$만 읽는다. 실제 결과 시점, 실제 구간 길이, 교전 이후의 참여자, 이후 상태, 최종 승패는 결코 입력이 아니다. 규모 계급(코호트 소속)은 예측 시점 \emph{이후}의 참여로 계산되는 사후 속성이므로 모델 입력이나 배치(어느 교전에 걸 것인가)의 근거로 쓰지 않는다.

라벨 쪽의 누출은 두 평가기 경로가 막는다. TRAIN 교전의 라벨을 만드는 $\Vfold{k}$의 전처리, 가중치, 조기 종료 보류, 보정기 어디에도 fold $k$의 경기 정보가 들어가지 않으며, 그 사례의 $\ppre$와 양 끝점에는 같은 fold 제외 평가기가 쓰인다. 학습 라벨용 평가기와 평가 역할용 동결 번들의 연결은 manifest로 기록된다. 이 규칙은 입력의 시간적 이용 가능성과 자기 경기의 최종 결과 이용을 점검하는 것이지, 라벨의 독립적 타당성이나 인과 식별을 보증하지 않는다.

외부 코호트의 평가는 점수 전용이다. $\Vhat$, $q$, 기준선, 전처리, 보정기 어느 것도 재적합하지 않으며, 사례 키나 특징 순서가 맞지 않으면 실패로 닫힌다(fail-closed; 특징 순서 sha16이 봉인되어 있다). 소규모 외부 코호트는 전이 성공 주장으로 합산하지 않는다.
